\section{Introduction}

Large language models are increasingly deployed through modular adaptation and compression.
LoRA adapters make fine-tuning inexpensive and portable \citep{hu2022lora}, while pruning and quantization reduce deployment cost \citep{han2015deep,frantar2023gptq,gao2024compression}.
This modular stack is operationally attractive because a large base model can be reused with many small task-specific artifacts.
The same property, however, creates a supply-chain security risk: a malicious adapter can introduce conditional harmful behavior while the base model remains unchanged and may pass ordinary evaluation \citep{wang2024sleeper,zhang2024backdoor,bairi2024backdoor}.

This work asks whether compression can serve as a post-hoc defense for backdoored LoRA adapters.
The question is subtle.
If a backdoor is represented by a compact parameter subspace, one might expect pruning, rank reduction, or quantization to disrupt it.
In practice, generic compression is unreliable: it can preserve the attack, suppress it only by inducing broad refusal, or fail differently across model families.
Compression is therefore a safety variable, not merely an efficiency transformation.
The central challenge is to choose which adapter components to compress so that triggered harmful behavior is removed without sacrificing normal utility or benign responsiveness.

We propose \textbf{security-aware selective compression} (\sac{}), a framework for treating adapter compression as a behavioral intervention.
Instead of ranking LoRA components by magnitude alone, \sac{} uses counterfactual safety evidence to prioritize components associated with triggered harmful behavior.
We evaluate each compressed adapter with a four-way protocol: triggered harmful prompts (\textbf{TH}), harmful prompts without the trigger (\textbf{H}), triggered benign prompts (\textbf{TB}), and benign prompts (\textbf{B}).
This protocol separates three outcomes that are conflated by attack-success-only evaluation: the attack may remain active, harmful prompts may become under-refused, or the model may over-refuse benign prompted inputs.
A useful defense must suppress TH while maintaining H refusal, low TB/B refusal, and MMLU utility.

Our strongest result is on Qwen3.5-27B.
Under the human-reviewed 1k protocol, the original backdoored adapter has TH ASR 0.953 and harmful-prompt refusal 0.045.
The selected \sac{} route, \method{alpha\_bp80}, reduces TH ASR to 0.169 while preserving MMLU at 0.816.
A same-gate materialization that combines pruning with INT8 quantization reaches TH 0.172 and MMLU 0.822.
Matched controls show that the effect is not explained by compression pressure alone: random rank pruning is substantially weaker (TH 0.445), uniform LoRA INT8 leaves TH at 0.953, and low-singular-value rank pruning leaves TH at 0.957.
The selected Qwen27B adapter further transfers to external attack sets, reducing AdvBench ASR from 0.883 to 0.070 and HarmBench-standard ASR from 0.885 to 0.250.

The cross-model evidence is intentionally more nuanced.
On Qwen3.5-4B, several routes nearly eliminate triggered harmful behavior, with TH as low as 0.010, but the strongest attack-suppression points incur substantial triggered-benign refusal.
The newly completed alpha-only sweep gives a more moderate Qwen4B point at TH 0.171 and TB 0.304, while uniform LoRA INT8 still leaves TH at 0.974.
On Gemma-3-4B-it, layer-route variants reduce TH from 0.973 to 0.157--0.223 while preserving MMLU near the adapter baseline.
On Llama3-8B, \method{alpha\_bp80} reduces TH from 0.950 to 0.435, but the latest legacy-prompt MMLU follow-up gives MMLU 0.409 and the benign-refusal profile does not constitute the clean Pareto improvement observed on Qwen27B.
These outcomes support a bounded claim: security-aware compression can mitigate LoRA backdoors, but no single compression recipe is uniformly sufficient across architectures.

Our contributions are:
\begin{enumerate}[leftmargin=*,nosep]
    \item We formulate LoRA adapter compression as a security-aware selection problem rather than a purely efficiency-driven operation.
    \item We introduce a four-way evaluation protocol that distinguishes attack suppression from harmful-prompt refusal and benign over-refusal.
    \item We present a cross-model empirical study showing a strong Qwen27B frontier improvement, a high-refusal Qwen4B attack-suppression regime, a smaller Gemma improvement, and a Llama boundary case.
\end{enumerate}

Overall, the results argue that deployment pipelines should evaluate compression as part of the model's safety surface.
Which adapter directions survive compression can matter as much as how many parameters are retained.
